\section{Two-stage A/B unblinding}
\label{sec:ab}

\subsection{The proposal}

Split the dataset into a fraction $f$ (sample A, exploration) and $1-f$ (sample B, confirmation).
Scan \emph{everything} in A with \emph{no} global correction, since what happens in A is selection
rather than inference. Pre-register every window with local $Z_A \ge Z_{\mathrm{cut}}$. Unblind
only those windows in B, and make the discovery claim on the B-only $p$-value, corrected for the
$k$ pre-registered windows.

The attraction is immediate. The trials factor in B is not modelled but \emph{counted}: it is the
number of lines in a frozen document. For a classical fixed-window search that is a modest
convenience. For a neural-network-driven scan, where ``how many independent looks did the network
take'' has no crisp answer, it is the difference between a defensible number and an argument.

\subsection{The cost}

Nothing is free, and here the trade is explicit. One avoids a logarithmic penalty by giving up
luminosity, which costs $\sqrt{f}$ on the significance. For a signal whose full-dataset
local significance would be $Z_{\mathrm{full}}$, the two stages fluctuate independently, so the
joint discovery probability is
\begin{equation}
  P_{\mathrm{disc}}
  \;=\;
  \Phi\!\left(\sqrt{f}\,Z_{\mathrm{full}} - Z_{\mathrm{cut}}\right)
  \cdot
  \Phi\!\left(\sqrt{1-f}\,Z_{\mathrm{full}} - Z_{B}^{\mathrm{req}}\right),
  \qquad
  Z_B^{\mathrm{req}} = \sqrt{25 + 2\ln k_{\mathrm{eff}}},
  \label{eq:abpower}
\end{equation}
with $k_{\mathrm{eff}} = N\,p_1(Z_{\mathrm{cut}}) + 1$ the expected number of pre-registered
windows. The reach is the $Z_{\mathrm{full}}$ at which $P_{\mathrm{disc}} = 50\,\%$.

The distinction between joint power and the naive median arithmetic matters. Requiring each stage
to succeed at its median gives $Z_{\mathrm{med}} = \max(Z_{\mathrm{cut}}/\sqrt{f},\,
Z_B^{\mathrm{req}}/\sqrt{1-f})$, which \emph{underestimates} the required signal wherever both
constraints bind, since two independent coin flips at 50\,\% each give 25\,\% jointly.
Equation~\eqref{eq:abpower} is the honest figure of merit, and toy Monte Carlo reproduces it to
$0.03\sigma$.

At this study's budget ($N = 6616$, single-stage reach $6.53\sigma$):

\begin{table}[htbp]
\centering
\caption{Signal strength required for a 50\,\% probability of discovery, as a function of the
splitting scheme, at $N = 6616$ independent looks.}
\label{tab:abcost}
\begin{tabular}{lrr}
\toprule
procedure & reach $Z_{\mathrm{full}}$ & cost \\
\midrule
single-stage full scan, exactly corrected & 6.53 & --- \\
optimised split ($f \approx 0.2$--$0.3$, $Z_{\mathrm{cut}} \approx 2$--$3$) & $\sim7.0$ & $+0.5\sigma$ \\
naive 50/50 split, $Z_{\mathrm{cut}} = 3$ & $\sim8.0$ & $+1.4\sigma$ \\
\bottomrule
\end{tabular}
\end{table}

\begin{figure}[htbp]
\centering
\includegraphics[width=0.9\textwidth]{ab_split_reach.png}
\caption{Discovery reach versus exploration fraction $f$, for several stage-1 selection thresholds.
Solid curves are 50\,\% joint power (Eq.~\eqref{eq:abpower}), thin dashed the naive median
arithmetic. The red line is the exactly corrected single-stage scan. The optimum is markedly asymmetric, with a small exploration sample and a large confirmation
sample, and the naive 50/50 split is a poor choice.}
\label{fig:abreach}
\end{figure}

\subsection{The look-elsewhere effect is conserved, not reduced}
\label{sec:abcons}

It is tempting to expect a trials count large enough to make splitting pay. There is none, and the
reason is an identity rather than a numerical accident. The cost $2\ln N$ factorises:
\begin{equation*}
  2\ln N \;=\; 2\ln\!\frac{N}{k} \;+\; 2\ln k,
\end{equation*}
where the first term is paid implicitly, as the stage-1 selection cut that reduces $N$ looks to
$k$ pre-registered windows, and the second explicitly in B. Splitting \emph{relocates} the penalty
rather than removing it. What remains is the two-coin penalty, the requirement that both
independent halves succeed, which runs from $+0.3$ to $+0.5\sigma$ between $N = 10$ and $N =
10^{10}$ and never crosses zero.

The number of unblinded windows barely matters: moving $k$ from 1 to 1000 changes the reach by
$0.15\sigma$, so $k$ can be chosen on practical grounds ($k \approx 10$).

\begin{figure}[htbp]
\centering
\includegraphics[width=0.92\textwidth]{ab_toys_power.png}
\caption{Toy Monte-Carlo validation: discovery power versus injected signal strength for the three
procedures, with the analytic median-reach values as vertical lines. The toys treat the budget's
$N$ effective looks as independent standard normals under background-only and recover the
full-dataset scan exactly as $z_{\mathrm{full}} = \sqrt{f}z_A + \sqrt{1-f}z_B$.}
\label{fig:abtoys}
\end{figure}

\subsection{Would swapping the halves recover it?}
\label{sec:abswap}

The natural objection to that identity is that half the luminosity sits idle at each stage. Why
not run \emph{both} directions --- explore in A and confirm in B, \emph{and} explore in B and
confirm in A --- and claim if either confirms? This is two-fold cross-validation of the
\emph{claim}, and it is a genuinely different procedure rather than a relabelling.

The trade is explicit. Both directions pre-register, so the confirmation windows double,
$k \to 2k$, costing $2\ln 2$ on the B bar: at $Z_{\mathrm{cut}} = 3$ and $N = 6616$,
$Z_B^{\mathrm{req}}$ moves from $5.64$ to $5.75$. What is bought is a second chance at the same
signal. Since $Z_B^{\mathrm{req}} > Z_{\mathrm{cut}}$, the two directions coincide only where
\emph{both} halves clear the claim bar, so the union power is
\begin{equation}
  P^{\mathrm{sym}}_{\mathrm{disc}}
  \;=\;
  \Phi^{A}_{\mathrm{cut}}\Phi^{B}_{\mathrm{req}}
  \;+\;
  \Phi^{B}_{\mathrm{cut}}\Phi^{A}_{\mathrm{req}}
  \;-\;
  \Phi^{A}_{\mathrm{req}}\Phi^{B}_{\mathrm{req}},
  \qquad
  \Phi^{A}_{x} = \Phi\!\left(\sqrt{f}\,Z_{\mathrm{full}} - Z_{x}\right),
  \label{eq:abswap}
\end{equation}
and $\Phi^{B}_{x}$ the same with $f \to 1-f$; toys reproduce it to $0.01\sigma$. At fixed
thresholds the gain over one direction is a factor
$2 - \Phi_{\mathrm{req}}/\Phi_{\mathrm{cut}}$, that is, at most two.

That gain is collected only when both directions have comparable power, which means $f = 1/2$ ---
and $f = 1/2$ is precisely the split the one-way analysis rejects, because confirmation is the
luminosity-hungry stage. The swap therefore does not beat the asymmetric scheme; it
\emph{rescues the symmetric one}.

\begin{table}[htbp]
\centering
\caption{One-way splitting against the symmetrised swap at $N = 6616$ (single-stage reach
$6.53$). The first two rows re-optimise $f$ at fixed $Z_{\mathrm{cut}}$; the last optimises
$(f, Z_{\mathrm{cut}})$ inside the practical design box $Z_{\mathrm{cut}} = 2$--$4.5$.}
\label{tab:abswap}
\begin{tabular}{lrrrr}
\toprule
 & \multicolumn{2}{c}{one-way} & \multicolumn{2}{c}{swapped} \\
\cmidrule(lr){2-3}\cmidrule(lr){4-5}
design & reach & cost & reach & cost \\
\midrule
$50/50$, $Z_{\mathrm{cut}} = 3$ & 7.98 & $+1.45$ & \textbf{7.38} & $\mathbf{+0.86}$ \\
$f = 0.30$, $Z_{\mathrm{cut}} = 3$ (recommended) & \textbf{7.09} & $\mathbf{+0.56}$ & 7.15 & $+0.63$ \\
best design in the practical box & \textbf{7.01} & $\mathbf{+0.49}$ & 7.03 & $+0.51$ \\
\bottomrule
\end{tabular}
\end{table}

Applied on top of the recommended working point the swap is actively harmful: the second
direction would have to confirm a $5\sigma$-global result on $30\%$ of the data, so it almost
never fires, and one pays the $2\ln 2$ for a dead direction. Optimised freely inside the
practical box the two schemes are a wash, $+0.49\sigma$ against $+0.51\sigma$, and the gap stays
within $0.11\sigma$ over nine decades of $N$, changing sign with the design box rather than
crossing over at any particular trials count.

Why neither can win is clearest geometrically. In the $(z_A, z_B)$ plane the Neyman--Pearson
acceptance region for a mean shift is the half-plane
$\sqrt{f}\,z_A + \sqrt{1-f}\,z_B > \sqrt{25 + 2\ln N}$ --- which is exactly the single-stage scan
on the recombined dataset. One-way splitting approximates that half-plane with an $L$-shaped
corner, the swap with a two-step staircase. More steps track the line more closely, which is why
the swap repairs the $50/50$ split, but every step multiplies the pre-registered windows.
Conservation of the look-elsewhere effect is that staircase-versus-half-plane inefficiency, and no
number of folds removes it.

The decisive objection, though, is not statistical. Of the three things the split buys
(Section~\ref{sec:abwhy}) the swap keeps two and destroys the third. The trials factor is still
counted and still auditable --- $2k$ is as countable as $k$ --- but procedural blindness is gone:
after the first direction the analyst has seen B, so the second direction's ``confirmation'' on A
is made by someone who already knows both halves. It is valid only if the entire two-directional
procedure is frozen and executed in a single automated pass, and a review committee is entitled to
discount it even then. The swap also does nothing about the correlated-systematics mode of
Section~\ref{sec:abblind}, which for a data-directed analysis is the leading one.

None of this argues against symmetrising the \emph{training}, which is a separate question taken
up in Section~\ref{sec:abnoise}: holding the background model out of the sample it scores is nearly
free and attacks the sculpting mode directly. Symmetrising the \emph{claim} is the thing that
costs, and on this budget it buys nothing.

\subsection{When the split is nonetheless the right choice}
\label{sec:abwhy}

The split wins if and only if the single-stage trials factor one would have to \emph{defend}
exceeds $R^\star \approx 13$ times the true $N$. For a classical fixed-window scan that is
implausible, and the corrected single stage wins outright. For a network-driven hunt a factor of
13 in forced conservatism is entirely imaginable, and that, rather than the logarithm, is the
quantitative case for two-stage unblinding. It buys three things:

\begin{itemize}[leftmargin=1.4em,itemsep=2pt]
  \item an \textbf{exactly countable} trials factor in B,
  \item \textbf{auditability}: the correction is the number of lines in a document frozen before
        unblinding, and
  \item \textbf{procedural blindness}: the analysis cannot be tuned on the confirmation sample.
\end{itemize}

\subsection{The failure mode splitting does \emph{not} address}
\label{sec:abblind}

A/B splitting attacks the \emph{statistical} look-elsewhere effect only. A spurious bump caused by
correlated systematics, such as background mis-modelling or a network that has learnt to sculpt a
peak, appears identically in both halves and will happily ``confirm''. For a data-directed
analysis this is not a corner case but the leading failure mode. The split protects against
fluctuations, not against the analysis fooling itself. That statement is exact but incomplete:
between a pure fluctuation and a fully correlated systematic lies everything an imperfect
estimator does, and Section~\ref{sec:abnoise} prices the split over that range.

\subsection{What changes if the significance estimator is imperfect}
\label{sec:abnoise}

Everything above assumes a perfectly calibrated estimator: every look returns $Z \sim N(0,1)$
under background, so the only thing to control is the look-elsewhere effect, and splitting can
only lose. A network-driven scan does not have that estimator. It has one that occasionally
returns an outlier --- a fit that failed to converge on this subsample, a background model that
has absorbed a fluctuation, a window where the shape is simply wrong. Repricing the split against
that estimator changes the verdict, and the reason it changes is worth more than the number.

The only property of an outlier that an A/B split can see is whether it \emph{repeats in the other
half}. That suggests three classes, distinguished by nothing but their correlation across the
split:

\begin{table}[htbp]
\centering
\caption{Classes of estimator defect, by how they correlate between the two halves. Only the last
column matters to a split.}
\label{tab:abdefect}
\begin{tabular}{lp{4.5cm}p{5.4cm}}
\toprule
class & physical origin & behaviour across the split \\
\midrule
\textsc{bias} & mismodelled background or detector artefact in that window &
  the pull is $\delta B/\sqrt{B} = \delta\sqrt{wB}$, so it grows with luminosity exactly like a
  signal and is \textbf{identical} in both halves \\
\textsc{var} & the look's uncertainty is underestimated by a factor $s$ &
  the miscalibration is shared, but what it amplifies is each half's \emph{independent}
  fluctuation \\
\textsc{glitch} & a per-\emph{run} failure: non-convergence on this subsample, a network
  artefact, a corrupted input &
  redrawn \textbf{independently} every time the estimator is run \\
\bottomrule
\end{tabular}
\end{table}

The mechanism is a one-line geometric statement, and Fig.~\ref{fig:abmech} is the whole section in
a picture. In the $(Z_A, Z_B)$ plane a coherent pull $b$ places a look at
$(\sqrt{f}\,b,\,\sqrt{1-f}\,b)$ --- on the locus $Z_B = \sqrt{(1-f)/f}\,Z_A$, which is precisely
where a real signal of the same strength sits. An incoherent one places it near the $Z_B \approx 0$
axis. The split is a cut that separates the axis from the locus, and by construction cannot
separate the locus from itself. It follows immediately that a coherent defect is not merely hard
for the split to reject: it is \emph{exactly as confirmable as a signal}, which is the content of
Section~\ref{sec:abblind} restated.

\begin{figure}[htbp]
\centering
\includegraphics[width=\textwidth]{ab_outliers_mechanism.png}
\caption{\textbf{Left:} looks that pass the stage-1 cut, in the plane of the two half-dataset
significances. Real signals and coherent mismodellings populate the same locus; per-run artefacts
sit on the $Z_B \approx 0$ axis and die at the claim bar. \textbf{Right:} the resulting
spurious-claim rate. The two defects are \emph{identical} to a single-stage scan --- one curve, by
construction --- and the split separates them by changing the scaling from $\epsilon$ to
$\epsilon^{2}$ for the incoherent one alone.}
\label{fig:abmech}
\end{figure}

The same two failure modes on an actual falling spectrum are shown in Fig.~\ref{fig:abspec}: a
background fit that undershoots by $\sim10\,\%$ in the A half alone produces a $4.1\sigma$
excess that is duly pre-registered and then collapses to $2.6\sigma$ in B, while a background
genuinely mismodelled by $8.8\,\%$ gives $4.2\sigma$ in A and $6.9\sigma$ in B and clears the claim
bar.

\begin{figure}[htbp]
\centering
\includegraphics[width=\textwidth]{ab_outliers_spectrum.png}
\caption{Toy spectra scored the way an analysis would score them, at the recommended working point.
\textbf{Left:} the A-half background model is locally wrong, so the excess exists only in the
estimator and stage B removes it. \textbf{Right:} the background model is wrong for both halves, so
the excess is real in the data and confirms. Only the first of these is a fluctuation in the sense
the split can act on.}
\label{fig:abspec}
\end{figure}

The comparison cannot be made on raw false-alarm rates, because the single stage can always buy
robustness by raising its bar. The honest figure of merit is therefore \emph{matched}: raise the
single-stage threshold until it is equally likely to make a spurious claim, and only then compare
reach. Table~\ref{tab:abnoise} does this at the recommended working point, with each defect
occurring at rate $\epsilon = 10^{-3}$ and magnitude $\sim2\sigma$.

\begin{table}[htbp]
\centering
\caption{Matched-robustness comparison at $f = 0.30$, $Z_{\mathrm{cut}} = 3$, $N = 6616$, for a
defect rate $\epsilon = 10^{-3}$ ($\approx 1$ artefact above $3\sigma$ per scan). A positive gain
means the split reaches a \emph{weaker} signal than a single-stage scan that is equally safe.}
\label{tab:abnoise}
\begin{tabular}{lrrrr}
\toprule
estimator & suppression & split reach & matched bar & gain \\
\midrule
perfect (the assumption above) & --- & 7.09 & 6.53 & $-0.56\sigma$ \\
\textsc{bias} & $1.9\times$ & 7.08 & 6.95 & $-0.14\sigma$ \\
\textsc{var} ($s = 2$) & $3.5\times$ & 7.10 & 7.20 & $+0.10\sigma$ \\
\textsc{glitch} & $199\times$ & 7.11 & \textbf{9.61} & $\mathbf{+2.50\sigma}$ \\
\bottomrule
\end{tabular}
\end{table}

Against a per-run glitch the split suppresses the spurious-claim rate by two orders of magnitude,
because the artefact must recur \emph{independently} in B above $5.6$: the rate goes as
$\epsilon^2$ rather than $\epsilon$. A single-stage scan would have to demand $Z \ge 9.6$ to be as
safe. Against a coherent bias the split suppresses by $1.9\times$, but that is not discrimination
--- it is merely the split's own worse reach hurting fakes along with signal --- and the single
stage buys the same factor for $+0.42\sigma$, which is cheaper than the split costs. Section
\ref{sec:abblind} is the $\rho = 1$ limit of this table.

\subsubsection*{Coherence is the deciding variable}

Interpolating between the two extremes with a coherent fraction $\rho$ (Fig.~\ref{fig:abnoise}),
the break-even sits at $\rho^\star \approx 0.80$ and moves by less than $0.02$ over two decades in
the overall defect rate. The rule is therefore unusually portable:

\begin{quote}
\textbf{If fewer than about four in five of the estimator's outliers repeat in the other half, the
split pays for itself.}
\end{quote}

Under a plausible cocktail --- coherent bias, bad fits and glitches each at $10^{-3}$, together
$\approx 3$ artefacts above $3\sigma$ per scan --- the net is $+0.29\sigma$ in the split's favour,
against $-0.56\sigma$ for a perfect estimator. And $\rho$ is measurable \emph{before} unblinding:
split the sidebands or MC-only pseudo-data, and count what fraction of the $3\sigma$ excesses
recurs.

\begin{figure}[htbp]
\centering
\includegraphics[width=0.92\textwidth]{ab_split_outliers.png}
\caption{Matched-robustness advantage of the split as a function of the coherent fraction $\rho$
of the estimator's outliers, at three defect rates. The dotted line is the pure cost of splitting
with a perfect estimator. The break-even is essentially rate-independent, so the criterion is a
property of the estimator's failure modes rather than of how often they occur.}
\label{fig:abnoise}
\end{figure}

\subsubsection*{Where the training arrangement enters}

The sharpest case is a background model \emph{fit to data}, which is the data-directed situation.
Fitting it on the same events one then tests for excesses is overfitting, and overfitting is a
\textsc{glitch}-type defect tied to the specific events in that half. A single pass cannot escape
it: there is no held-out half. Three arrangements are available to a split, and they are not
equivalent:

\begin{enumerate}[leftmargin=1.6em,itemsep=3pt]
  \item \textbf{self-contained} --- train on A and score A, train on B and score B. Both stages
        overfit, but independently, so the claim is $\epsilon^2$-suppressed. Bookkeeping exact.
        Worth $+2.50\sigma$.
  \item \textbf{full cross} --- train on B and score A, train on A and score B. No overfitting
        anywhere, worth $+5.67\sigma$, \emph{but} the stage-A selection is now a function of B
        through the trained model, so the pre-registration is no longer independent of the
        confirmation sample and the $k$-coin $p$-value is no longer exact.
  \item \textbf{A-trained throughout} --- train on A, score A to select, then apply that same
        frozen model to B. The confirmation $Z$ is produced by an estimator that has never seen a
        single B event, so overfitting is removed \emph{by construction} rather than beaten down
        statistically; and nothing about B enters the pre-registration, so the bookkeeping stays
        exact. Worth $+5.64\sigma$.
\end{enumerate}

Arrangement 3 dominates: it recovers all but $0.03\sigma$ of the full-cross gain while keeping the
blindness that made the scheme worth adopting. The reason it can is that the overfitting it leaves
in place --- in stage A --- only inflates the candidate list ($k$: $8.9 \to 10.1$), and $k$ enters
the bar as $2\ln k$. \textbf{Contaminating the exploration stage is logarithmically cheap;
contaminating the confirmation stage is what costs.}

\subsubsection*{Two caveats on this result}

Part of the \textsc{glitch} gain is available without splitting. If the artefact is purely
seed- or initialisation-dependent, re-running the estimator with different seeds on the
\emph{same} data catches it at no cost in reach. What a seed jackknife cannot decorrelate is a
defect tied to a specific \emph{data} fluctuation --- overfitting, or a fit driven by the actual
events in that window. Only a data split touches those, which is exactly the data-directed failure
mode, so the claim is narrower than the raw $+2.5\sigma$ but lands on the case at hand.

Second, and more important: after the split the residual spurious rate is essentially
\emph{entirely} coherent. The glitch and variance components contribute $0.9\,\%$ and $0.1\,\%$ of
what still gets through. The split does not make the analysis safe --- it changes what one has to
worry about, from ``is my estimator glitchy?'' to ``is my background model wrong in that window?''.
That residual has to be controlled by other means (validation regions, spurious-signal tests,
alternative background parameterisations), because the confirmation stage will endorse a
mismodelled background as enthusiastically as a real signal.

The overall repricing is then a change of argument rather than of arithmetic.
Section~\ref{sec:abcons} stands exactly as written: with a perfect estimator, splitting relocates
the look-elsewhere effect and cannot reduce it. But that was answering the wrong question for a
data-directed search. The split is not a look-elsewhere device --- it is an
\textbf{estimator-robustness device}, and priced as one, $0.5\sigma$ is cheap.

\subsection{Recommended working point}

If the scheme is adopted, the design that follows from the above is:

\begin{itemize}[leftmargin=1.4em,itemsep=2pt]
  \item exploration fraction $f \approx 0.25$--$0.35$, stage-1 threshold $Z_{\mathrm{cut}} = 3$
        (giving $\approx 9$ pre-registered windows at this budget);
  \item pre-registered $\pm2\sigma_M$ windows, with the outcome ladder written down before
        unblinding;
  \item the background model \textbf{trained on A alone} and applied frozen to B
        (arrangement 3 of Section~\ref{sec:abnoise}), so that the confirmation significance comes
        from an estimator that has never seen a B event, without leaking B into the
        pre-registration;
  \item the claim quoted as a \textbf{global} significance over the pre-registered windows:
        evidence at $Z_B^{\mathrm{global}} \ge 3$ (local $Z_B \ge 4.0$), discovery at
        $Z_B^{\mathrm{global}} \ge 5$ (local $5.6$). The false-evidence rate per opening is then
        $1.35\times10^{-3}$ by construction, whatever number of windows A flags;
  \item B alone for the claim $p$-value, A$+$B for the measurement.
\end{itemize}
